% Contenido del póster: Pampurini, Pezzola y Quaranta (2024), Finance Research Letters. DOI 10.1016/j.frl.2024.105519
% Elaborado con apoyo de IA a partir del resumen oficial del artículo (OpenAlex), 2026-11-12.
\def\DOI{10.1016/j.frl.2024.105519}
\def\TituloEs{Modelos de negocio crediticios y eficiencia de las fintech}
\def\TituloOriginal{Lending business models and FinTechs efficiency}
\def\Autores{Pampurini, Pezzola y Quaranta}
\def\Anio{2024}
\def\Revista{Finance Research Letters}

\def\Problema{Las fintech de crédito se multiplican, pero se sabe poco sobre qué decisiones de gestión determinan su eficiencia.}
\def\Objetivo{Analizar qué aspectos de gestión son los principales determinantes de la eficiencia de todas las fintech italianas dedicadas al crédito.}
\def\Metodo{Medición de la eficiencia de las fintech de crédito italianas entre 2020 y 2022 con análisis envolvente de datos estocástico (SDEA) y estudio de sus determinantes.}
\def\Resultados{\begin{itemize}
  \item Los principales determinantes de la eficiencia son la rentabilidad sobre activos (ROA) y el ratio costo/ingreso.
  \item Controlar el nivel de riesgo del negocio y los costos es clave para los gestores de estas empresas.
\end{itemize}}
\def\Conclusion{Los resultados sirven a las fintech, a otros actores financieros interesados en fusiones y adquisiciones, y a reguladores y supervisores que buscan reglas homogéneas en el sector del crédito.}
\def\Relevancia{El estudio justifica usar el ROA para evaluar a las sociedades administradoras peruanas. En el artículo del autor, solo la plataforma líder obtuvo utilidades, con un ROA anualizado de \cifra{roaLider2025}~\% en 2025; las demás cerraron con pérdidas en \cifra{aniosOtrasPerdida} de sus \cifra{aniosOtras} años con información.}
\def\PalabrasClave{fintech, eficiencia, SDEA, rentabilidad sobre activos, ratio costo/ingreso}
% --- Ficha del estudio (herramientas/posteres/ficha_articulos.py) ---
\def\Ambito{Italia, 2020-2022 (fintech de crédito)}
\def\Enfoque{Análisis envolvente de datos estocástico (SDEA)}
\def\FiguraPeru{4}
% --- fin de la ficha ---
